
\vspace{-0.3cm}
\section{Related Work}
\vspace{-0.3cm}

\paragraph{Structural Analysis of Neural Networks.}
Prior work has studied neural-network structure through architecture graphs, weight spectra, and representation similarity. Architecture-level analyses relate graph properties such as clustering and path length to predictive performance, while weight-space methods based on random matrix theory connect spectral structure to training quality, generalization, and pruning sensitivity~\cite{you2020graph, martin2021implicit}. Representation-similarity methods such as SVCCA and CKA compare hidden states across layers and models, but do not explicitly model interactions among model components~\cite{raghu2017svcca, kornblith2019similarity, liu2025spectral}. For transformer-based models, prior studies show that attention heads are heterogeneous: some encode positional, syntactic, coreferential, or induction-like behavior, while many are redundant or weakly used~\cite{clark2019does, voita2019analyzing, michel2019sixteen}. LLM pruning methods exploit this heterogeneity through weight- or layer-level importance estimates, including SparseGPT, Wanda, and ATP~\cite{frantar2023sparsegpt, sun2023simple, huang2025determining}. In contrast, our work constructs activation-defined functional graphs over attention heads, quantifies their global small-world organization, and derives local head- and layer-level roles from this graph structure.

\paragraph{Functional Connectivity and Small-world Organization.}
Functional connectivity measures statistical dependence among distributed neural activities~\cite{friston1994functional}. Graph-theoretic work characterizes small-world networks by high clustering and short path lengths~\cite{watts1998collective}, or by high local and global efficiency~\cite{latora2001efficient}. In neuroscience, small-world organization is viewed as a balance between segregated specialization and integrated communication in brain networks~\cite{bassett2006small, bullmore2009complex}. Related machine-learning work uses activation-based graphs, activation correlations, and topological summaries to study neural-network organization and generalization~\cite{horta2021extracting, ballester2024predicting, ballester2023topological}, while transformer representation analyses reveal clustered structure across layers~\cite{phang2021fine}. Building on this perspective, we use SWI to test whether attention-head activations in LLMs form small-world functional organization associated with reasoning capability.